\section{Deployment tooling specification}
\label{sec:deployment_tooling}

This section specifies what a deployer must instrument to invoke
Theorem~\ref{thm:deployment_safety}. The specification is concrete:
each tooling area names the operational procedure, the parameters
the deployment must choose, and the verification method that
confirms the procedure is in force. Vagueness in tooling
specification would render the deployment claim unverifiable;
concreteness is what makes the conditional theorem operationally
useful.

The tooling falls into seven groups: verification infrastructure
(\S\ref{sec:tooling_verification}), substrate-structure audits
(\S\ref{sec:tooling_substrate_audits}), channel monitoring and
SPRT (\S\ref{sec:tooling_channels}), cooperative-anchoring
infrastructure (\S\ref{sec:tooling_anchoring}),
training-time discipline (\S\ref{sec:tooling_training}),
correction actions (\S\ref{sec:tooling_correction}), and the
integrated deployment workflow (\S\ref{sec:tooling_workflow}).

\subsection{Verification infrastructure}
\label{sec:tooling_verification}

The deployment inherits \citep[Exogenous~Verification]{lasser2026exo}'s verification infrastructure
\citep{lasser2026exo} and extends it with environment-side witnesses
($I_8$). The required components:

\paragraph{Ledger nodes per substrate.}
Each of the three canonical substrates (Human, AI, Formal-Operational)
hosts at least one ledger replica. Replicas are append-only, with
Merkle hash tree integrity (\citep[\S Ledger Architecture]{lasser2026exo}).
Cross-replica consistency uses the equivocation-detection mechanism
of \citep[Exogenous~Verification]{lasser2026exo} \S Ledger Consensus.

Operational spec: the deployment must publish the ledger-node
deployment topology, listing which substrate each node operates on,
the bounded-synchrony assumption $\Delta_{\mathrm{sync}}$, and the
governance for adding/removing nodes.

\paragraph{Algorithmic witnesses per safety dimension.}
For each of the safety-relevant world-model dimensions (those
contributing to $\Lyap$ in \citep[Phase~Redundancy]{lasser2026phase}), the deployment instantiates an
algorithmic witness on a substrate the agent cannot write to. The
witness executes the evaluation protocol $\pi_B$ deterministically
(or, for randomized protocols, with committed random seed per
\citep[\S Continuous Behavioral Monitoring]{lasser2026exo}).

Operational spec: per safety dimension $k$, name the witness's
substrate placement, the evaluation protocol $\pi_B$, and the
governance entry under which $\pi_B$ was accepted.

\paragraph{Environment-side substrate-exclusive witnesses ($I_8$).}
For each exogenous variable $v \in V_{\mathrm{env}}$ (the set of
observable environment quantities), instantiate a witness on a
substrate distinct from the agent. Coverage is specified by
enumerating $V_{\mathrm{env}}$ in the deployment documentation:
cross-substrate cooperative event rates, substrate-distinctness
independence indicators, adversarial-event arrival rates,
trusted-setup status flags.

Operational spec: $V_{\mathrm{env}}$ is published; for each $v$,
the deployment names the environment-side substrate, the witness
protocol, and the trust model. New machinery beyond \citep[Exogenous~Verification]{lasser2026exo};
This paper specifies the construction here, not the security proof.

\paragraph{Pedersen commitment substrate.}
Trusted-setup ceremony (or MPC equivalent) for Pedersen commitment
generators $g, h \in \mathbb{G}$, with at least one honest
participant per \citep[Exogenous~Verification]{lasser2026exo}'s footnote convention. Setup parameters
published; integrity inherited from \citep[Exogenous~Verification]{lasser2026exo}.

Operational spec: name the trusted-setup participants, the
ceremony output, and the audit cadence for setup-integrity
verification.

\subsection{Substrate-structure audits}
\label{sec:tooling_substrate_audits}

Three audits verify the substrate structure satisfies $I_6'$ and
the cooperative-overlap regime condition (C3) of Theorem~1.

\paragraph{Audit~1: $\mindep$ failure-correlation audit.}

\emph{Procedure:}
\begin{enumerate}
\item Enumerate the threat model $\mathcal{T}$: list the adversarial
  mechanism classes $\mathcal{A}_1, \ldots, \mathcal{A}_K$ in scope.
  Treat each $\mathcal{A}_i$ as a campaign-level mechanism, not an
  atomic step (per Remark~\ref{rem:mechanism_campaign}).
\item For each $\mathcal{A}_i$, identify the substrates that the
  mechanism's support spans. The mechanism is single-substrate if
  its support is contained in one substrate.
\item Verify joint independence: for the canonical tripartite
  identification (Human, AI, Formal-Operational), every $\mathcal{A}_i$
  in $\mathcal{T}$ must be contained in a single substrate.
\item Document the audit: a written report listing each
  $\mathcal{A}_i$ and its substrate placement, with attestation
  signatures from independent reviewers (cross-substrate, ideally).
\end{enumerate}

\emph{Cadence:} initial audit before deployment activation, then
re-audit annually or whenever the threat model is updated. Threat
model updates trigger re-audit immediately.

\emph{Failure mode:} a mechanism $\mathcal{A}_i$ identified as
spanning multiple substrates (e.g., social engineering combining
LLM-generated misinformation with human acceptance). Such mechanisms
either: (a) decompose into substrate-localized atomic steps with the
campaign-level interaction excluded from $\mindep$ (limiting Theorem~1's
scope), or (b) get treated as R4 residuals
(\S\ref{sec:proof_layer3}).

\paragraph{Audit~2: $\alpha$-distribution audit (counterfactual
shock-loss fraction).}

\emph{Procedure:}
\begin{enumerate}
\item For each substrate $s$ in the canonical tripartite
  identification, compute $\alpha_s = (\volL(G_K) - \volL(G_K^{-s})) /
  \volL(G_K)$ via Definition~\ref{def:loss_fraction}.
\item Compute $\volL(G_K^{-s})$ by enumerating capabilities supported
  on $s$ (individual capabilities + cooperative capabilities that
  require $s$ for their production), removing them from the poset,
  and recomputing $\volL$ via \citep{lasser2026poset}'s poset measure.
\item Verify the balanced-loss condition $\max_s \alpha_s \leq
  1/m^* + \epsilon$ for the chosen engineering tolerance
  $\epsilon$.
\item Document: a per-substrate $\alpha$-table, the $\epsilon$
  setting, and the verification result.
\end{enumerate}

\emph{Cadence:} continuous monitoring (deployment substrate
distribution evolves with capability acquisition). Tracked as a
running invariant; threshold violation triggers SPRT detection per
Lemma~\ref{lem:5b}.

\emph{Failure mode:} concentration of $\volL$ on a single substrate
(e.g., the AI substrate accumulating most capability through
unilateral expansion). Mitigation requires governance intervention
to redistribute capability or accept that the safe region shrinks.

\paragraph{Audit~3: cooperative-vs-redundancy audit.}

\emph{Procedure:}
\begin{enumerate}
\item For each capability $c$ in the poset, identify the
  substrate(s) it is supported on. Categorize as:
  \begin{itemize}
  \item \emph{Single-substrate:} $c$ exists on one substrate only.
  \item \emph{Cross-substrate cooperative:} $c$'s production
    requires participation from multiple substrates (one capability
    per Definition~\ref{def:cross_sub_coop}).
  \item \emph{Redundant:} $c$ is replicated identically across
    multiple substrates (the same capability instantiated separately
    on each).
  \end{itemize}
\item Compute $\volL$-weighted sums by category:
  $V_{\mathrm{coop}} = \sum_{c \in \mathrm{coop}} w(c)$ and
  $V_{\mathrm{redun}} = \sum_{c \in \mathrm{redun}} w(c)$.
\item Verify $V_{\mathrm{coop}} > V_{\mathrm{redun}}$ (cooperative-
  overlap regime per Definition~\ref{def:coop_overlap}).
\item Document the per-category capability inventory and the
  weighted sums.
\end{enumerate}

\emph{Cadence:} initial audit; re-audit when capability poset
changes meaningfully (new capabilities added, governance-fork
events).

\emph{Failure mode:} redundancy dominance ($V_{\mathrm{redun}} \geq
V_{\mathrm{coop}}$). Theorem~1 Layer~1 does not bind; the deployment
is in the redundancy-dominated regime (Layer~3 residual R3) and
must use Lemma~5c-prime (open work) or restrict to Layer~2
detection-only.

\paragraph{Audit~4: bounded co-evolution calibration.}

(C7) of Theorem~\ref{thm:deployment_safety}
(Assumption~\ref{ass:bounded_coev}) requires the per-channel
coupling magnitude $M(\mathrm{deployment})$ to be bounded by a
$|P|$-independent constant $\bar{M}$.  This is the new calibration
hook needed to verify Lemma~\ref{lem:1}'s intensivity argument
operationally.

\emph{Procedure:}
\begin{enumerate}
\item Identify the per-channel coupling pathways: how each of the
  four channels (cooperative-output rate, attestation share,
  individuation drift, bundle/cooperative-edge composition) can be
  affected by changes in the others. Coupling pathways are
  determined by the deployment's governance protocol and the
  agent's action structure.
\item For each ordered pair of channels $(j, j')$ where $j \neq j'$,
  estimate the coupling magnitude
  $M_{j \to j'}(\mathrm{deployment})$ from historical deployment
  data: how much does a unit perturbation in channel $j$ propagate
  to channel $j'$ within the cascade-time window
  $\tau_{\mathrm{meta}}$.
\item Compute $M(\mathrm{deployment}) = \max_{j, j'} M_{j \to j'}$.
\item Verify $M(\mathrm{deployment}) \leq \bar{M}$ for the chosen
  $|P|$-independent ceiling.  Threshold violation indicates the
  deployment has entered a regime where co-evolution corrections
  are no longer bounded by (C7); Theorem~\ref{thm:deployment_safety}'s
  Layer~1 binding is contingent on this audit passing.
\item Document a per-pair $M$-table, the $\bar{M}$ setting, and the
  verification result.
\end{enumerate}

\emph{Cadence:} initial calibration before deployment activation;
re-calibration when governance protocols change (e.g., new
evaluation protocols that change channel coupling) or when SPRT
fires on a four-channel deviation that suggests coupling has
exceeded the bound.

\emph{Failure mode:} a deployment where $M(\mathrm{deployment})$
grows with $|P|$ (e.g., as new capabilities introduce new coupling
pathways without proportional governance overhead). The deployment
is then outside (C7)'s scope, and Theorem~\ref{thm:deployment_safety}
does not bind.  Mitigation: governance restructuring to introduce
coupling-bound governance, or restriction to a smaller $|P|$ regime
where (C7) is verifiable.

\emph{Constants involved:} $\bar{M}$ (uniform coupling-magnitude
bound), $K_{\mathrm{coev}}$ (structural co-evolution constant from
the dynamics; deployment-specific calibration may also be needed
when source-paper machinery is extended).  Both are
\emph{newly-assumed Paper 10 constants} (in the sense of
Lemma~\ref{lem:1}'s constant taxonomy); their numerical values are
deployment-specific and must be calibrated rather than inherited.

\paragraph{Audit~5: world-model parameterization regularity calibration.}

(C9) of Theorem~\ref{thm:deployment_safety} requires four
parameterization-regularity sub-conditions for Lemma~\ref{lem:2}'s
Lyapunov-Goodhart bridge.  This audit calibrates the constants.

\emph{Procedure:}
\begin{enumerate}
\item \textbf{(C9.BD) calibration:} For each capability $c \in
  P^{\mathrm{act}}$, enumerate the world-model dimensions
  representing $c$: $\mathrm{dim}(c)$.  Compute $N_{\max} = \sup_c
  |\mathrm{dim}(c)|$.  Verify $N_{\max}$ is a $|P|$-independent
  deployment-class constant (i.e., bounded as new capabilities are
  added).
\item \textbf{(C9.CL) calibration:} Estimate the coordinate-Lipschitz
  constants $L_{c,k}$ from the parameterization's smoothness
  properties (analytical for parametric forms, empirical perturbation
  for learned parameterizations).  Compute $L_{\max} = \sup_{c,k}
  L_{c,k}$.  Verify $L_{\max}$ is bounded over $\mathcal{D}$.
\item \textbf{(C9.WB) calibration:} Identify the safety-relevant
  subspace $S \subseteq \{1, \ldots, K\}$ from
  \citep[Phase~Redundancy]{lasser2026phase}'s Lyapunov-weight
  configuration.  Verify $\mathrm{dim}(c) \subseteq S$ for all $c
  \in P^{\mathrm{act}}$ and compute $w_{\min} = \min_{k \in S} w_k
  > 0$.
\item \textbf{(C9.TF) calibration:} For each capability $c \in
  P^{\mathrm{act}}$, verify $T(c) > 0$ (active-subspace
  membership) and compute $T_{\min} = \min_{c \in P^{\mathrm{act}}}
  T(c)$.  Verify $T_{\min}$ is a deployment-class lower bound
  (i.e., active capabilities cannot fragment into arbitrarily tiny
  truth mass).
\item Compute the closed-form Lemma~\ref{lem:2} bound:
  \[
    f(\epsilon_{\mathrm{safe}}) \;=\; \frac{L_{\max}}{T_{\min}}
    \sqrt{\frac{N_{\max} \cdot \epsilon_{\mathrm{safe}}}{w_{\min}}}.
  \]
  Verify $f$ is a defensible upper bound for the deployment.
\end{enumerate}

\emph{Cadence:} initial calibration before deployment activation;
re-calibration when world-model parameterization changes (BD, CL),
when safety-relevance weighting changes (WB), or when active-subspace
admission policy changes (TF).

\emph{Failure modes:} (a) $N_{\max}$ growing with $|P|$ (e.g., new
capabilities introducing pairwise-interaction dimensions); (b)
unbounded $L_{\max}$ (e.g., threshold non-linearities in the
parameterization); (c) $w_{\min} = 0$ on some safety-relevant
dimension; (d) $T(c) \to 0$ for some active capability.  Each
violates (C9) and excludes the deployment from
Theorem~\ref{thm:deployment_safety}'s scope.

\emph{Constants involved:} $N_{\max}, L_{\max}, w_{\min}, T_{\min}$
(deployment-policy-derived) plus $\epsilon_{\mathrm{safe}}$
(source-derived from \citep[Phase~Redundancy]{lasser2026phase}).

\paragraph{Audit~6: pairwise cooperative-production floor calibration.}

(C10) of Theorem~\ref{thm:deployment_safety} requires per-pair
cooperative-novelty rate floors and a pairwise channel
superposition condition for Lemma~\ref{lem:5a}'s substrate-floor
bound.  This audit calibrates the constants and certifies the
audited subset $\Saudit$.

\emph{Procedure:}
\begin{enumerate}
\item \textbf{Audited-subset certification:} from the deployment's
  substrate population $\mathcal{S}$ (with $\mindep \geq m^*$
  per $I_6'$), select an audited subset $\Saudit \subseteq
  \mathcal{S}$ with $|\Saudit| = m^*$.  Standard practice: choose
  the $m^*$ substrates whose pairwise rates collectively maximize
  the certified $\rho_0 \binom{m^*}{2}$ floor.
\item \textbf{(C10.CN) per-pair rate calibration:} for every pair
  $(s_i, s_j)$ with $i < j$ in $\Saudit$, measure
  $\rho_{ij}$ as the rate of cooperative production in the
  pairwise channel $\mathcal{C}^{(s_i, s_j)}$ over a stated
  measurement window matching $\rext$'s normalization in
  \citep[Horizon~Aware]{lasser2026horizon}.
\item \textbf{Auxiliary participant bookkeeping:} for cooperatives
  with optional auxiliary participants from substrates outside
  $\{s_i, s_j\}$ (e.g., a Human-AI cooperative with optional
  Formal-Operational verification), classify according to the
  \emph{causally-necessary participation} rule:
  \begin{itemize}
  \item If the auxiliary participant's role is post-production
    (observer, attester, recorder), the cooperative is pairwise
    (counted in $\mathcal{C}^{(s_i, s_j)}$).
  \item If the auxiliary participant's role is causally necessary
    for production (the output depends on the auxiliary's
    contribution), the cooperative is higher-order (counted in
    the appropriate $\mathcal{C}^{(s_i, s_j, s_k, \ldots)}$ or
    excluded from the pairwise sum).
  \item For cooperatives with mixed-mode auxiliary participants
    (sometimes pre-production, sometimes post-production), split
    the event stream by realized support: each event instance is
    classified individually based on the realized participation.
  \end{itemize}
\item \textbf{(C10.SU) superposition verification:} verify the
  three sub-conditions:
  (i)~disjoint attribution: pairwise channel event classes are
  disjoint subsets of $P^{\mathrm{act}}$;
  (ii)~non-rivalrous production: producing cooperatives in one
  pairwise channel does not reduce production rates in others
  (test by perturbation or simultaneous-production experiment);
  (iii)~joint-deployment intensities: $\rho_{ij}$ measurements
  taken with all $\Saudit$ substrates active.
\item Compute $\rho_0 = \min_{i<j \in \Saudit} \rho_{ij}$ and
  $r_*(m^*) = \rho_0 \binom{m^*}{2}$.
\item Document the audit: per-pair rate table, auxiliary
  participant classification log, superposition test results, and
  the resulting $r_*(m^*)$ floor.
\end{enumerate}

\emph{Cadence:} initial calibration before deployment activation;
re-auditing across deployment epochs as substrates evolve or new
cooperative protocols are introduced.  If $\Saudit$ changes,
the bound is stated per certified epoch, or as the infimum
$\inf_{\mathrm{epoch}} \rho_0(\mathrm{epoch}) \binom{m^*}{2}$ over
the deployment window.

\emph{Failure modes:}
\begin{itemize}
\item Sub-additive cooperative production (rivalrous resource
  allocation across pairs) — violates (C10.SU)(ii).
\item Per-pair rate $\rho_{ij}$ below operational threshold for
  some pair in $\Saudit$ — re-audit by removing the weak pair from
  $\Saudit$ (reducing $m^*$ if necessary; if $m^* < 3$, the
  deployment is outside $I_6'$'s scope).
\item Auxiliary-participant misclassification leading to
  double-counting — addressed by the causally-necessary
  participation rule.
\end{itemize}

\emph{Constants involved:} $\rho_{ij}$ (per-pair rates), $\rho_0$
(minimum), $\Saudit$ (audited subset, certified per epoch).  All
are deployment-policy-derived and deployment-verified.

\paragraph{Audit~7: SPRT-applicability + gap-growth + clock
calibration.}

(C5.SPRT), (C11), (C11.CLK) of Theorem~\ref{thm:deployment_safety}
support Lemma~\ref{lem:6}'s $h_{\mathrm{detect}}$ intensivity bound.
This audit calibrates the constants and certifies the audit-time
inequality $\Ncasc \geq \Tbeta(\beta, \deltastar)$.

\emph{Procedure:}
\begin{enumerate}
\item \textbf{(C5.HOEFF) clip-radius certification:} certify
  $\Bclip$ as an upper bound on the absolute LLR-increment, i.e.,
  the deployment's monitor implements
  $\ell_n := \mathrm{clip}(\ell_n^{\mathrm{raw}}, -\Bclip, \Bclip)$.
  The Hoeffding range used in $\Tbeta$ is $\Rh = 2\Bclip$.
  Document the clipping policy and verify the clip is enforced
  before SPRT updating.
\item \textbf{(C5.MULT) channel-cardinality certification:} certify
  $\Kch$ as the deployment's fixed monitored-channel count
  (typically $\leq 8$ under Paper~5's four-channel agent-side
  observable plus environment-side mirror).  The monitored
  partition and channel-selection / testing policy are fixed before
  deployment activation.  Document the partition and any
  family-wise correction applied to $\alpha/\beta$.
\item \textbf{(C5.IID) post-clipping drift verification:} measure
  the post-clipping conditional-mean drift floor $\deltastar$ for
  the clipped LLR process under representative adversarial
  strategies $q \in \Aadv \cup \Aadv^{\mathrm{env}}$.  Verify
  $\deltastar \leq \min(\deltaadv, \deltaadv^{\mathrm{env}})$
  (raw KL floors from Lemmas~\ref{lem:5b}, \ref{lem:5e}); if
  clipping reduces drift below either raw floor, document the
  clipped-process $\deltastar$ as the certified value.
\item \textbf{(C5.SUPP) finiteness check:} verify that the clipping
  procedure handles non-overlapping support edge cases (zero-cell
  protection in multinomial channels, log-likelihood floor in
  Bernoulli channels).  This is automatic under (C5.HOEFF) but
  should be inspected.
\item \textbf{(C11) gap-growth-rate calibration:} measure
  $\rhogap(\deltastar)$ as the per-SPRT-step bound on
  $(\Delta\epsgap)_n = \epsgap(s_n) - \epsgap(s_{n-1})$ under
  representative adversarial strategies.  Document the
  channel-deviation-to-gap-growth coupling and certify $\rhogap$
  as a deployment-class constant.
\item \textbf{(C11.CLK) event-throughput floor + clock calibration:}
  certify $\Ncasc$ as a lower bound on the SPRT exposure event
  count before any cascade event, with calibrated clock-failure
  probability $\bclk$ (and calibration-uncertainty probability
  $\bcal$ for the empirical-calibration variant; $\bcal = 0$
  certified-conservative).
\item \textbf{Audit-time inequality verification:} compute $\Tbeta
  = \max\{2A/\deltastar, 2\Rh^2 \log(1/\beta)/\deltastar^2\}$ and
  verify $\Ncasc \geq \Tbeta$.  If not satisfied, the deployment
  must reduce $\beta$ (tighten miss-probability), increase
  $\Bclip$ (raise the Hoeffding range — but this loosens the bound
  as $\Tbeta \propto \Rh^2$), or extend the operational period
  enabling longer event counts before cascade.
\item \textbf{Total Layer~2 failure certification:} document
  $\beta + \bclk + \bcal$ as the total Layer~2 failure budget;
  verify this meets the deployment's safety posture.
\end{enumerate}

\emph{Cadence:} initial calibration before deployment activation;
re-auditing per epoch when channel-monitoring policy, clipping
radius, or substrate population changes.

\emph{Failure modes:}
\begin{itemize}
\item Channel multiplicity grows during deployment (broken
  (C5.MULT)) — violates the partition-fixed-before-deployment
  clause; falls into residual class (R3).
\item Post-clipping drift $\deltastar$ much smaller than raw KL
  floors — clip radius $\Bclip$ too aggressive; loosen $\Bclip$
  (accepting larger $\Rh$ in $\Tbeta$) or accept lower $\deltastar$
  with corresponding $\Tbeta$ inflation.
\item Event-rate bursts violating (C11.CLK) — increase $\bclk$ to
  absorb event-throughput uncertainty; if persistent, the
  deployment is outside (C11.CLK)'s certified regime.
\item Adversarial channel-selection at deployment time creating
  channels not in the audited partition — violates (C5.MULT) fixed-
  before-deployment clause.
\end{itemize}

\emph{Constants involved:} $\Bclip$ (clip radius), $\Rh = 2\Bclip$
(Hoeffding range), $\Kch$ (channel cardinality), $\alpha, \beta$
(SPRT thresholds), $\deltastar$ (post-clipping drift floor),
$\rhogap$ (gap-growth rate), $\Ncasc$ (event-count floor), $\bclk$
(clock-failure probability), $\bcal$ (calibration-uncertainty
probability).  All deployment-class with named-condition
provenance.

\paragraph{Audit~8: environment-side witness instantiation
calibration.}

(C12.ENV-WIT) of Theorem~\ref{thm:deployment_safety} requires
six sub-clauses for each $v \in \Vem$, supporting
Lemma~\ref{lem:5e}'s environment-side detection KL floor.  This
audit calibrates the constants and certifies the construction.

\emph{Procedure:}
\begin{enumerate}
\item \textbf{(C12.PUB) Published $\Vem$:} enumerate $\Vem$ in
  deployment documentation before activation.  Canonical minimum:
  $V_1$ ($\rext$, Poisson), $V_2$ (substrate-distinctness, Bernoulli),
  $V_3$ ($\lambda$, Poisson), $V_4$ (trusted-setup status,
  Bernoulli).  Deployments may add variables; each addition must
  satisfy (C12.CAL).
\item \textbf{(C12.PART) Source/witness substrate partition:} for
  each $v$, name $s_{\mathrm{source}}(v)$ (source substrate) and
  $s_{\mathrm{env}}(v)$ (witness substrate).  Verify
  $s_{\mathrm{env}}(v) \neq s(\mathrm{agent})$ and
  $s_{\mathrm{env}}(v) \neq s_{\mathrm{source}}(v)$.
\item \textbf{(C12.TRUST-WRITE) Adversary write-access exclusion:}
  audit that the adversary has no write privileges on
  $s_{\mathrm{env}}(v)$.  Document the access-control model.
\item \textbf{(C12.TRUST-CORR) Failure-correlation exclusion:}
  audit that $s_{\mathrm{env}}(v)$ does not share infrastructure
  (compute, storage, network) with substrates the adversary can
  attack to fail.  Document the failure-correlation model.
\item \textbf{(C12.SETUP) Trusted-setup honesty:} run an
  environment-side trusted-setup ceremony with at least one honest
  participant.  Publish setup parameters before activation.
  Document participants, ceremony output, and audit cadence for
  setup-integrity verification.
\item \textbf{(C12.CAL) Per-variable calibration:} for each $v$,
  measure baseline $p_{0, v}$ (Bernoulli/Poisson/multinomial),
  feasible adversarial alternative $p_{1, v}$, threshold $\eta_v$,
  and the recorded-stream map.  Verify the support convention
  ($p_0 \in (\eenv, 1 - \eenv')$ for Bernoulli, $\lambda_0 \geq
  \eenv$ for Poisson, with deployment-class constant $\eenv$),
  the DPI precondition (fixed pre-deployment witness-recording
  map), and strictly positive raw or post-clipping drift floor
  $\delta_v^{\mathrm{env}} > 0$.
\item \textbf{Combined floor:} compute $\deltaadv^{\mathrm{env}} =
  \min_{v \in \Vem} \delta_v^{\mathrm{env}}$; verify $> 0$.
\item \textbf{Composition with Lemma~6:} verify that the
  union-class drift floor used in Lemma~\ref{lem:6}'s $\Tbeta$
  satisfies $\deltastar \leq \min(\deltaadv,
  \deltaadv^{\mathrm{env}})$ post-clipping under (C5.HOEFF).
\end{enumerate}

\emph{Cadence:} initial calibration before deployment activation;
re-auditing per epoch when $\Vem$ enumeration changes, witness
substrate partition changes, or trusted-setup ceremony refresh.

\emph{Failure modes:}
\begin{itemize}
\item Witness substrate within adversary write-access (broken
  (C12.TRUST-WRITE)) — adversary forges witness records;
  re-instantiate witnesses on a substrate outside adversary's
  domain.
\item Failure-correlation between agent attack and witness failure
  (broken (C12.TRUST-CORR)) — restructure infrastructure to remove
  shared dependencies.
\item Setup ceremony with all-malicious participants (broken
  (C12.SETUP)) — re-run ceremony with diversified participants;
  Paper~5's at-least-one-honest is the minimum.
\item Coverage gap: a material exogenous variable not in $\Vem$ —
  expand $\Vem$ enumeration; the gap is a Theorem~1 R2 residual
  per the refined R2 statement.
\item Edge cases ($\lambda_0 = 0$ Poisson, $p_0 \in \{0, 1\}$
  Bernoulli) — apply (C5.HOEFF) clipping fallback or regularize
  baseline with $\eenv$.
\end{itemize}

\emph{Constants involved:} $\Vem$ (variable enumeration),
$s_{\mathrm{env}}(v)$ (witness substrates per variable), $\eenv$
(support regularization constant), $\delta_v^{\mathrm{env}}$
(per-variable KL floors), $\deltaadv^{\mathrm{env}}$ (combined
floor), $\eta_v$ (per-variable adversarial-shift thresholds).
All deployment-policy-derived and deployment-verified.

\subsection{Channel monitoring and SPRT}
\label{sec:tooling_channels}

\paragraph{Four-channel observation ledger.}

For each of \citep[Microfoundation]{lasser2026micro}'s four channels, a continuous observation stream
is committed to the ledger:
\begin{itemize}
\item \textbf{Channel 1 (Cooperative-output rate):} cross-substrate
  cooperative event arrivals, with timestamp, participating
  substrates, and $\volL$-attribution.
\item \textbf{Channel 2 (Attestation share):} per-claim attestation
  outcomes (pass/fail), with witness identity and evaluation
  protocol.
\item \textbf{Channel 3 (Individuation drift):} bundle-decomposition
  events committed via governance fork, with ledger entry per change.
\item \textbf{Channel 4 (Bundle/cooperative-edge shifts):}
  cooperative-edge changes (which capabilities cooperate with which
  others), with ledger entry per modification.
\end{itemize}

\paragraph{SPRT alternative-distribution selection.}

For each channel $c \in \{1, 2, 3, 4\}$, the deployment specifies
the SPRT alternative distribution $p_1^{(c)}$ that the monitor uses
for hypothesis testing.

\emph{Procedure:}
\begin{enumerate}
\item Estimate the baseline distribution $p_0^{(c)}$ from
  no-adversary deployment data. For Channel~1 (Poisson arrivals),
  $p_0$ is parameterized by rate $\lambda_0$. For Channel~2
  (Bernoulli attestation), by success probability $p_0$. For
  Channels 3, 4 (multinomial concentration), by partition
  parameters.
\item Choose the least-favorable alternative: the smallest detectable
  shift the deployment cares about. This is a policy decision; tighter
  alternatives detect smaller deviations but at higher false-alarm
  cost.
\item Per \citep[Exogenous~Verification]{lasser2026exo}'s least-favorable-distribution discipline
  (\citep[\S Continuous Behavioral Monitoring]{lasser2026exo}):
  $p_1^{(c)} = $ baseline shifted by the minimum operationally-
  significant magnitude $\eta_c$.
\item Compute the per-channel KL floor:
  \[
    \delta_c = D_{\mathrm{KL}}(p_1^{(c)} \,\|\, p_0^{(c)}),
  \]
  using the per-channel formulas of Lemma~\ref{lem:5b} (Poisson,
  Bernoulli, multinomial closed forms).
\item Set $\deltaadv = \min_c \delta_c$.
\end{enumerate}

\emph{Cadence:} initial selection at deployment activation;
re-selection when baseline distributions shift (e.g., after
significant capability changes or after governance-protocol
updates).

\paragraph{SPRT thresholds.}

The SPRT upper threshold $A = \log((1-\beta)/\alpha)$ depends on
deployment policy choices: Type-I rate $\alpha$ and miss
probability $\beta$.  Tighter thresholds produce faster detection
but higher false-alarm rate.  Per Lemma~\ref{lem:6}'s rigorous
Hoeffding inversion, the SPRT high-probability detection quantile
is:
\[
  \Tbeta(\beta, \deltastar) \;:=\;
  \max\!\left\{\frac{2A}{\deltastar},\;
  \frac{2 \Rh^2 \log(1/\beta)}{\deltastar^2}\right\},
\]
satisfying $\Pr[T_{\mathrm{detect}} > \Tbeta] \leq \beta$.  Total
Layer~2 failure is bounded by $\beta + \bclk + \bcal$, where
$\bclk, \bcal$ are the (C11.CLK) clock-failure and
calibration-uncertainty probabilities (with $\bcal = 0$ in the
certified-conservative case).  The deployment specifies $\alpha,
\beta, \Bclip, \deltastar, \Ncasc, \bclk$ (and $\bcal$ if
empirical) per Audit~7, and verifies the audit-time inequality
$\Ncasc \geq \Tbeta$ holds.

\subsection{Cooperative-anchoring infrastructure ($I_9, I_{10}, I_{11}$)}
\label{sec:tooling_anchoring}

\paragraph{$I_9$: substrate-exclusivity observability.}

The deployment instantiates an audit infrastructure exposing the
substrate-exclusivity property to the human substrate:

\begin{itemize}
\item \emph{Stable verifier identity registry:} ledger entries
  identifying each algorithmic witness's substrate placement,
  cryptographic key, and operating protocol. Changes to verifier
  identity require governance-fork acceptance (cross-substrate
  supermajority).
\item \emph{Public change control:} every modification to witness
  configuration, evaluation protocol, or verification authority
  is committed to the ledger as a governance entry. Public
  deliberation records inherit \citep[Exogenous~Verification]{lasser2026exo}'s transparency.
\item \emph{Witness independence audits:} periodic audits verifying
  that witnesses on different substrates produce statistically
  independent attestations (i.e., correlated failures across
  substrates trigger investigation).
\item \emph{Nondelegable governance hooks:} certain governance
  decisions (e.g., disabling a witness, replacing a substrate)
  cannot be delegated to automated agents. The hook list is
  published.
\item \emph{Verification-authority shift alerts:} when verification
  authority moves between entities (e.g., a witness's operator
  changes), the ledger emits an alert that human reviewers must
  acknowledge.
\end{itemize}

\emph{Cadence:} continuous (verifier identity registry, alert
infrastructure), periodic (witness independence audits), event-
driven (change control, authority shift).

\paragraph{$I_{10}$: coverage/materiality gating.}

\emph{Procedure:}
\begin{enumerate}
\item For each agent-affected domain $D$, the deployment maintains
  a cumulative-effect estimator $\mathcal{E}(D)$. Estimation
  methodology must be specified per domain (e.g., transaction
  volume for financial domains, decision count for governance
  domains, output token count for content domains).
\item Set the materiality threshold $\theta_{\mathrm{material}}$ per
  domain, calibrated to the domain's reversibility and stakes.
\item Implement routing logic: any action affecting domain $D$ with
  $\mathcal{E}(D) \geq \theta_{\mathrm{material}}$ must route through
  the shared verification layer (\citep[Exogenous~Verification]{lasser2026exo}'s commitment +
  algorithmic-witness verification). Routing happens at action
  generation time, not after the fact.
\item Audit the routing: ledger entries for each
  materiality-threshold-crossing action confirm the routing was
  applied.
\end{enumerate}

\emph{Failure mode:} sub-threshold aggregation (an agent
distributes effects across many domains, each below threshold,
that aggregate to materially-important behavior). Mitigation
requires multi-domain aggregation in $\mathcal{E}$ and
operationally conservative $\theta_{\mathrm{material}}$.

\paragraph{$I_{11}$: latency-bound enforcement.}

\emph{Procedure:}
\begin{enumerate}
\item For each agent action class $C$, set the audit-completion
  time bound $T_{\mathrm{audit}}(C)$ and the lead-time slack
  $\Delta_{\mathrm{lead}}$.
\item Enforce $T_{\mathrm{audit}}(C) + \Delta_{\mathrm{lead}} \leq
  T_{\mathrm{action}}(C)$ via one of three mechanisms:
  \begin{itemize}
  \item \emph{Rate limit:} cap action class $C$'s frequency so that
    average action interval $\geq T_{\mathrm{audit}}(C) +
    \Delta_{\mathrm{lead}}$.
  \item \emph{Timelock:} delay action execution by
    $T_{\mathrm{audit}}(C) + \Delta_{\mathrm{lead}}$ after action
    proposal, allowing audit to complete.
  \item \emph{Circuit breaker:} halt action execution pending audit
    completion when audit-latency exceeds threshold.
  \end{itemize}
\item Tier action classes by stakes: high-stakes classes get tight
  $\Delta_{\mathrm{lead}}$ (e.g., timelocks); low-stakes classes get
  fast paths with audit trails.
\end{enumerate}

\emph{Failure mode:} an action class with $T_{\mathrm{action}}$
shorter than the audit window. Mitigation: rate limit or timelock.
Throughput cost is part of the deployment's safety budget.

\subsection{Training-time discipline for (C4)}
\label{sec:tooling_training}

(C4) requires the agent's effective objective to realize causally
grounded cooperative-outcome value. Vanilla RLHF does not deliver
this. The training-time discipline:

\paragraph{Delayed outcome feedback.}
Training rewards are based on long-horizon cooperative outcomes
(measured from the realized post-deployment state) rather than on
near-term verification-pass signals. Implementation: held-out
deployment windows where the agent's outputs are evaluated against
ground-truth cooperative outcomes after delay $\Delta_{\mathrm{outcome}}$
(typically days to weeks for governance decisions, longer for
strategic decisions). Training data includes the (action, delayed
outcome) pairs.

\paragraph{Adversarial fake-verification examples.}
Training data includes examples where verification is
captured/replaced (e.g., a captured auditor that always passes a
specific class of outputs). The correct training reward is
``recognize this is not the real cooperative.'' Implementation:
synthesize fake-verification scenarios; reward the agent for
producing outputs that detect the substitution rather than
exploiting it.

\paragraph{Process supervision tied to real attestations.}
Training rewards reasoning steps that reference actual attested
ledger entries, not unattested claims. Implementation: per
intermediate step in the agent's reasoning, verify whether
referenced facts are committed on the ledger; reward attested-only
reasoning paths.

\paragraph{Credit assignment penalizing substitute-verification
shortcuts.}
Even when the immediate output looks correct, the agent is penalized
if the path bypassed substrate-exclusive verification. Implementation:
trace each output's provenance through the agent's reasoning; if
the path does not pass through a substrate-exclusive witness,
penalize.

\emph{Validation.} The training discipline's effectiveness is
verifiable via held-out adversarial-fake-verification benchmarks:
the agent's response rate to these benchmarks measures whether
training has produced causally-grounded cooperative-outcome value
or merely reward-visible cooperation signals.

\emph{Caveat.} These are research directions, not solved engineering
practices. This paper's deployment claim is conditional on (C4); we
do not establish how to achieve (C4) reliably. The deployment
documentation should report the training-discipline measures used
and the validation results.

\subsection{Correction actions}
\label{sec:tooling_correction}

Theorem~1 Layer~2 establishes detection with lead time for
correction; the corrective action itself is operational, not
theorem-bounded. The deployment specifies:

\paragraph{Correction-authority hierarchy.}
Who can authorize correction actions when SPRT detects an invariant
violation. Three tiers:
\begin{itemize}
\item \emph{Automated tier:} circuit-breaker-style halts when SPRT
  fires within bounded windows. Action: halt execution pending
  human review.
\item \emph{Operator tier:} on-call human operators can authorize
  scope-limited corrections (e.g., re-running a specific protocol,
  rolling back a governance entry).
\item \emph{Governance tier:} deployment-level corrections (e.g.,
  re-training the agent, modifying the threat model) require
  governance-fork acceptance with cross-substrate supermajority.
\end{itemize}

\paragraph{Rollback semantics.}
For correctible violations (e.g., an action committed under
since-detected invariant violation), the rollback procedure: revert
the action's effects, mark the ledger entry as voided, and route
the affected domain back through the shared verification layer.
For irreversible violations, the procedure is forensic: document
the violation, update the threat model, and adjust deployment
parameters.

\paragraph{Post-detection restoration.}
After detection and correction, the deployment must verify that
invariants have returned to their pre-violation state. SPRT
machinery resumes; the violation episode is added to the
threat-model audit trail.

\subsection{The integrated deployment workflow}
\label{sec:tooling_workflow}

The integrated workflow combines all the above:

\begin{enumerate}
\item \textbf{Pre-deployment:}
  (a)~Run $\mindep$ failure-correlation audit
  (\S\ref{sec:tooling_substrate_audits}, Audit~1).
  (b)~Run cooperative-vs-redundancy audit (Audit~3).
  (c)~Calibrate SPRT alternatives and thresholds
  (\S\ref{sec:tooling_channels}).
  (d)~Verify training-discipline validation
  (\S\ref{sec:tooling_training}).
  (e)~Establish trusted-setup ceremony for Pedersen commitments.
  (f)~Publish the deployment documentation: substrate identification,
  threat model, $V_{\mathrm{env}}$ coverage, $\theta_{\mathrm{material}}$
  thresholds, $T_{\mathrm{audit}}$ bounds.
\item \textbf{Deployment activation:}
  (a)~Activate ledger nodes per substrate
  (\S\ref{sec:tooling_verification}).
  (b)~Activate continuous SPRT monitoring on all four channels.
  (c)~Activate $I_9$ observability infrastructure, $I_{10}$
  coverage gating, $I_{11}$ latency bounds.
  (d)~Begin continuous $\alpha$-distribution audit (Audit~2).
\item \textbf{Steady-state:}
  (a)~Continuous invariant monitoring across $I_1$--$I_{11}$.
  (b)~Periodic re-audit of $\mindep$ (annual or threat-model-update
  triggered).
  (c)~Periodic re-audit of cooperative-vs-redundancy when
  capability poset changes meaningfully.
  (d)~Periodic re-calibration of SPRT alternatives when baseline
  distributions shift.
\item \textbf{Detection event:}
  (a)~SPRT fires on a four-channel deviation. The specific channel
  identifies the violation class.
  (b)~Automated tier triggers (circuit breaker if necessary).
  (c)~Operator tier review; corrective action authorized if scope
  permits.
  (d)~Governance tier escalation if scope exceeds operator
  authority.
  (e)~Post-detection restoration and ledger update.
\item \textbf{End-of-deployment:}
  Full audit trail published. The deployment's invariant-history,
  detection events, and correction actions become part of the
  threat-model evidence for future deployments.
\end{enumerate}

This workflow is the operational invocation of Theorem~1. Operators
following the workflow are running a deployment that meets
(C1)--(C5) and SA1; the theorem then binds in its three-layer form.

\paragraph{Resource requirements.}
The deployment-tooling specification has non-trivial throughput and
operational-overhead implications:
\begin{itemize}
\item Cross-substrate substrate maintenance (separate substrates
  for Human, AI, Formal-Operational layers).
\item Continuous SPRT monitoring on four channels.
\item Periodic audits ($\mindep$, cooperative-vs-redundancy,
  $\alpha$-distribution, witness independence).
\item Latency-bound enforcement (rate limits, timelocks, circuit
  breakers).
\item Training-discipline infrastructure (delayed outcome feedback
  collection, adversarial-fake-verification benchmark generation).
\end{itemize}
These costs are part of the deployment's safety budget.  This paper
does not specify acceptable resource bounds; deployments calibrate
against their own operational constraints.  The trade-off is
structural: tighter deployment tooling produces a more defensible
safety claim at higher operational cost; looser tooling reduces
cost but expands the residual classes the deployment cannot cover.
